\section{Performance}

\noindent
It was found that natural language processing on Javadoc is useful for classifying security-sensitive methods. A fine-tuned UniXcoder model was able to classify security-sensitive methods based on their Javadoc alone with an F1-score of $0.65$. \\

\noindent
However, it was also found that models trained on Javadocs exclusively did not perform as well as models using bimodal data (i.e., Javadoc \emph{and} code). This is likely a result of the specific pre-trained encoders being used. Since CodeBERT and UniXcoder were trained primarily on bimodal inputs, it is expected that these encoders perform best on input that more closely resembles the pre-training data. \\

\noindent
In almost all cases, the more sophisticated models (fine-tuned Transformer encoders with neural network classifiers) performed better than the more naive approaches (SVM classifiers, tf-idf feature generation, and pre-trained encoders without fine-tuning). In the case of models with bimodal input, the most simplistic model (an SVM classifier with tf-idf for feature extraction) had an F1-score of $0.51$ on the test dataset, whereas the best performing model (CodeBERT with mean pooling using a neural network classifier) achieved an F1-score of $0.68$. \\

\noindent
The performance of the Transformer encoders without fine-tuning was surprising. UniXcoder without fine-tuning, when paired with an SVM classifier, had an accuracy of $41\%$, which is considerably better than randomly guessing (i.e., $25\%$). In fact, the only model of this category that performed worse than randomly guessing was CodeBERT using CLS pooling which had an accuracy of $24\%$.\\

\noindent
Surprisingly, there were many cases of a significant difference in performance when using mean pooling versus CLS pooling. In general, it was found that UniXcoder performed better with mean pooling, and CodeBERT performed better with CLS pooling. At this stage, the cause of this discrepancy is unclear. \\

\noindent
Our models' performances fall short of Sas et al, SuSi, and SWAN, which achieved F1-scores of $86.4\%$, $91.9\%$, and $83\%$, respectively. However, this does not take away from the success of this thesis project. This thesis has shown that Javadoc contains information useful for security-related tasks, namely security-sensitive method classification. The approach presented in this thesis is not mutually exclusive to existing implementations and, rather than replace, could enhance current approaches. Furthermore, our models could be used by security analysts to generate a preliminary taint analysis configuration which could then be manually refined. This would likely result in a significant reduction in the time and effort required to produce a taint analysis configuration compared to creating one entirely from scratch. 

\section{Limitations and Improvements}
\label{limitations_and_improvements}

\noindent
There are some limitations to our approach. As is common with most machine learning projects, there was not enough high-quality labelled data to train and validate the models optimally. Ideally, the dataset should be larger and come from a greater diversity of projects. This would hopefully allow the models to generalise more effectively and perform better.\\

\noindent
Due to this lack of data, many examples needed to be manually labelled for this project. This is not ideal and undoubtedly lowered the quality of the dataset. A third party should review the dataset to ensure there are no mislabelled examples and that it is high-quality.\\

\noindent
Another limitation is in regards to how overloaded methods were handled. For this thesis, if a method was considered security-sensitive, all overloaded methods sharing its name were also considered security-sensitive. It is unclear if this assumption is always valid and is, therefore, a potential source of inaccuracy. \\